%───────────────────────────────────────────────────────────────────────────────
\chapter{Entornos teoremáticos}
\label{ch:entornos}
%───────────────────────────────────────────────────────────────────────────────

Los entornos teoremáticos son la interfaz más visible del template. Cada
uno tiene un rol epistémico preciso: no son simplemente cajas de color
distintos, sino que comunican al lector la naturaleza lógica del contenido.
Una \textbf{Definición} establece vocabulario. Un \textbf{Teorema} es una
verdad demostrada. Un \textbf{Ejemplo} instancia lo abstracto. Confundirlos
degrada el argumento matemático tanto como una demostración incorrecta.

Este capítulo documenta los ocho entornos disponibles, su sintaxis, su rol
semántico, y los casos donde cada uno es el correcto. Al final se cubren los
entornos especiales y el manejo de figuras y tablas dentro de cajas.

%───────────────────────────────────────────────────────────────────────────────
\section{Arquitectura de los entornos}
\label{sec:entornos-arquitectura}
%───────────────────────────────────────────────────────────────────────────────

\begin{definicion}{Generador de entornos}
\label{def:generador}
Todos los entornos teoremáticos se crean con la macro
\texttt{\textbackslash crearEntorno} definida en
\texttt{environments/generator.tex}~\cite{tcolorbox2024b}:
\begin{lstlisting}[style=latex, frame=none, numbers=none]
\crearEntorno{nombre}{color}{Titulo}
{conContador}{opcionesAdicionales}
\end{lstlisting}
El parámetro \texttt{color} es un alias semántico del dominio \texttt{env-*}
definido en la paleta. El fondo de la caja (\texttt{color-back}) se genera
automáticamente mediante \texttt{\textbackslash createback} al definir la paleta.
Ningún archivo de config toca colores \texttt{Ctp*} directamente.
\end{definicion}

La sintaxis de uso es uniforme para todos los entornos: un argumento
obligatorio con el título de la instancia.

\begin{lstlisting}[style=latex, frame=none, numbers=none]
\begin{nombre}{Título de esta instancia}
\label{tipo:etiqueta}   % opcional pero recomendado
Contenido del entorno...
\end{nombre}
\end{lstlisting}

Los entornos con contador incluyen automáticamente el número en el título
(\textit{Definición~1.1 --- Título}). Los entornos sin contador muestran
solo el título.

%───────────────────────────────────────────────────────────────────────────────
\section{Nivel 1: Entornos fundacionales}
\label{sec:entornos-nivel1}
%───────────────────────────────────────────────────────────────────────────────

Los entornos fundacionales establecen la estructura lógica del argumento.
Tienen contador para que puedan referenciarse con precisión.

\subsection{Definición}
\label{subsec:entorno-def}

Color: \texttt{env-def} $\to$ \textcolor{env-def}{\textbf{Blue}}.
Establece vocabulario formal. El contenido de una definición no se discute
ni se demuestra --- simplemente fija el significado de un término dentro
del sistema. Si el lector no conoce el término, la definición es su
contrato de lectura.

\begin{definicion}{Espacio métrico}
\label{def:espacio-metrico}
Un \textbf{espacio métrico} es un par $(\eqdom{X}, \eqfn{d})$ donde $\eqdom{X}$ es un conjunto
y $\eqfn{d} : \eqdom{X} \eqop{\times} \eqdom{X} \eqop{\to} \mathbb{R}_{\geq 0}$ es una \textit{métrica} que
satisface para todo $\eqk{x}, \eqk{y}, \eqk{z} \in \eqdom{X}$:
\begin{enumerate}[label=(\roman*), leftmargin=2em]
  \item $\eqfn{d}(\eqk{x},\eqk{y}) = 0 \iff \eqk{x} = \eqk{y}$ \hfill (identidad)
  \item $\eqfn{d}(\eqk{x},\eqk{y}) = \eqfn{d}(\eqk{y},\eqk{x})$ \hfill (simetría)
  \item $\eqfn{d}(\eqk{x},\eqk{z}) \leq \eqfn{d}(\eqk{x},\eqk{y}) + \eqfn{d}(\eqk{y},\eqk{z})$ \hfill (desigualdad triangular)
\end{enumerate}
\end{definicion}

\subsection{Teorema}
\label{subsec:entorno-thm}

Color: \texttt{env-thm} $\to$ \textcolor{env-thm}{\textbf{Mauve}}.
El color más prominente del sistema, reservado para el resultado más
importante de su categoría. Un teorema es una afirmación no trivial cuya
verdad se ha demostrado formalmente. La prominencia visual señala al lector
que aquí hay un resultado que merece atención sostenida.

\begin{teorema}{Completitud de los espacios $L^p$}
\label{thm:lp-completo}
Sea $(\eqdom{X}, \eqdom{\mathcal{M}}, \eqk{\mu})$ un espacio de medida y $1 \leq \eqk{p} < \infty$.
El espacio $\eqdom{L^p}(\eqdom{X}, \eqk{\mu})$ dotado de la norma
\[
  \|\equ{f}\|_{\eqk{p}} = \left(\eqop{\int}_{\eqdom{X}} |\equ{f}|^{\eqk{p}} \, \eqdm{d\mu}\right)^{1/\eqk{p}}
\]
es un espacio de Banach. En particular, $\eqdom{L^2}(\eqdom{X},\eqk{\mu})$ es un espacio de
Hilbert con producto interior $\langle \equ{f}, \equ{g} \rangle = \eqop{\int}_{\eqdom{X}} \equ{f}\,\overline{\equ{g}}\,\eqdm{d\mu}$.
\end{teorema}

\begin{observacion}{Teorema vs.\ Proposición}
Un \textbf{teorema} es el resultado principal de una línea de argumento ---
el punto al que todo lo demás conduce. Una \textbf{proposición} es un
resultado intermedio correcto y no trivial, pero subordinado. Si tienes
dudas, pregunta: ¿citarías este resultado de forma independiente, sin
contexto? Si sí, es un teorema. Si solo tiene sentido dentro del argumento
local, es una proposición.
\end{observacion}

\subsection{Proposición}
\label{subsec:entorno-prop}

Color: \texttt{env-prop} $\to$ \textcolor{env-prop}{\textbf{Green}}.
Resultado verdadero y demostrable, pero de alcance más local que un teorema.
El verde evoca una afirmación positiva --- algo que se sostiene.

\begin{proposicion}{Continuidad de la norma}
\label{prop:norma-continua}
Sea $(\eqdom{X}, \|\cdot\|)$ un espacio normado. La función $\|\cdot\| : \eqdom{X} \to \mathbb{R}$
es continua respecto a la topología inducida por la norma. Más aún, es
Lipschitz con constante $1$: para todo $\eqk{x}, \eqk{y} \in \eqdom{X}$,
\[
  \big|\|\eqk{x}\| - \|\eqk{y}\|\big| \leq \|\eqk{x} - \eqk{y}\|.
\]
\end{proposicion}

%───────────────────────────────────────────────────────────────────────────────
\section{Nivel 2: Entornos de apoyo}
\label{sec:entornos-nivel2}
%───────────────────────────────────────────────────────────────────────────────

Los entornos de apoyo complementan el argumento sin avanzarlo directamente.
No tienen la misma jerarquía lógica que los fundacionales, pero son
igualmente indispensables para la comprensión.

\subsection{Ejemplo}
\label{subsec:entorno-ej}

Color: \texttt{env-ex} $\to$ \textcolor{env-ex}{\textbf{Peach}}.
Instancia un concepto abstracto en un caso concreto. El tono cálido invita
a la lectura --- los ejemplos deben sentirse accesibles. Un buen ejemplo
no solo ilustra: muestra por qué la hipótesis general es necesaria (el
ejemplo borderline) o por qué el resultado no es trivial.

\begin{ejemplo}{$\ell^2$ como espacio de Hilbert}
\label{ex:l2-hilbert}
El espacio de sucesiones de cuadrado sumable
\[
  \eqdom{\ell^2} = \left\{(\equ{x}_n)_{n=1}^\infty \subset \mathbb{R} \;\middle|\;
  \eqop{\sum}_{n=1}^\infty \equ{x}_n^2 < \infty\right\}
\]
es un espacio de Hilbert~\cite{reed1980} con producto interior
$\langle \equ{x}, \equ{y} \rangle = \eqop{\sum}_{n=1}^\infty \equ{x}_n \equ{y}_n$.
La base ortonormal canónica $\{\eqk{e}_n\}_{n=1}^\infty$, donde $\eqk{e}_n$ tiene
un $1$ en la posición $n$ y ceros en el resto, es una base de Schauder
de $\ell^2$ pero \textit{no} es una base algebraica (Hamel): ningún
conjunto finito genera el espacio.
\end{ejemplo}

\subsection{Observación}
\label{subsec:entorno-obs}

Color: \texttt{env-obs} $\to$ \textcolor{env-obs}{\textbf{Overlay0}}.
El entorno más sutil del sistema --- casi gris. Una observación no avanza
el argumento principal ni instancia un concepto: llama la atención sobre
un detalle que el lector podría pasar por alto. Sirve para matices,
advertencias técnicas, notas de implementación, y condiciones que se
asumen implícitamente. El color discreto señala que puede omitirse en
una primera lectura sin perder el hilo.

\begin{observacion}{Convergencia débil vs. fuerte en $L^2$}
En espacios de Hilbert de dimensión infinita, la convergencia débil
no implica convergencia en norma. La sucesión $\eqk{e}_n \rightharpoonup 0$
débilmente en $\eqdom{\ell^2}$ (porque $\langle \eqk{e}_n, \equ{x} \rangle = \equ{x}_n \to 0$
para todo $\equ{x} \in \eqdom{\ell^2}$), pero $\|\eqk{e}_n\| = 1$ para todo $n$. Este fenómeno
no tiene análogo en dimensión finita y es fuente frecuente de errores en
argumentos de compacidad.
\end{observacion}

\subsection{Conexión física}
\label{subsec:entorno-phys}

Color: \texttt{env-phys} $\to$ \textcolor{env-phys}{\textbf{Sapphire}}.
Relaciona un resultado matemático abstracto con su interpretación o
motivación física. En documentos de física computacional o matemática
aplicada, este entorno ancla la abstracción en la intuición del dominio.
No es un ejemplo --- es un puente entre dos lenguajes.

\begin{fisicaconexion}{Completitud de $L^2$ y mecánica cuántica}
El \cref{thm:lp-completo} para $p=2$ es el fundamento del formalismo
de Hilbert en mecánica cuántica. El espacio de estados de una partícula
cuántica es $\eqdom{L^2}(\eqdom{\mathbb{R}^3})$: las funciones de onda $\equ{\psi}(\eqvec{\mathbf{r}})$
son elementos de $\eqdom{L^2}$, y la norma $\|\equ{\psi}\|_2^2 = \eqop{\int} |\equ{\psi}|^2 \, \eqdm{d^3r}$
es la probabilidad total (normalizada a 1). La completitud garantiza que
los límites de sucesiones de Cauchy de estados físicos son también estados
físicos --- sin ella, la evolución temporal $\equ{\psi}(t) = e^{-\eqcon{i}\eqop{\hat{H}}t/\eqcon{\hbar}}\eqk{\psi_0}$
podría salir del espacio.
\end{fisicaconexion}

%───────────────────────────────────────────────────────────────────────────────
\section{Nivel 3: Entornos especiales}
\label{sec:entornos-nivel3}
%───────────────────────────────────────────────────────────────────────────────

Los entornos especiales tienen propósitos concretos que los diferencian
de los niveles anteriores. Su uso debe ser deliberado.

\subsection{Estrella Polar}
\label{subsec:entorno-estrella}

Color: \texttt{env-star} $\to$ \textcolor{env-star}{\textbf{Yellow}}.
El amarillo llama la atención más que cualquier otro color del sistema.
Reservado para el resultado o idea que orienta toda una unidad temática ---
el punto hacia el que converge el argumento del capítulo o sección, como
una estrella polar orienta la navegación. No es simplemente un teorema
importante: es el \textit{por qué} existe el capítulo. Debe haber a lo
sumo uno por capítulo, y muchos capítulos no lo necesitan.

\begin{estrella}{Teorema espectral en espacios de Hilbert}
\label{star:espectral}
Sea $\eqop{A} : \eqdom{\mathcal{H}} \to \eqdom{\mathcal{H}}$ un operador autoadjunto acotado
en un espacio de Hilbert separable. Existe una medida espectral
$\eqfn{E} : \mathcal{B}(\mathbb{R}) \to \mathcal{L}(\eqdom{\mathcal{H}})$ tal que
\[
  \eqop{A} = \eqop{\int}_{\sigma(\eqop{A})} \equ{\lambda} \, \eqdm{dE}(\equ{\lambda}),
\]
donde la integral converge en la topología de operadores. Todo cálculo
funcional $\eqfn{f}(\eqop{A})$ se define de forma consistente como
$\eqfn{f}(\eqop{A}) = \eqop{\int} \eqfn{f}(\equ{\lambda})\,\eqdm{dE}(\equ{\lambda})$.
\end{estrella}

\subsection{Ejercicio Computacional}
\label{subsec:entorno-comp}

Color: \texttt{env-comp} $\to$ \textcolor{env-comp}{\textbf{Pink}}.
Para tareas que requieren implementación, simulación o verificación
numérica. El tono activo señala que el lector debe hacer algo, no solo
leer. A diferencia de los ejercicios matemáticos ordinarios (que usan
\texttt{\textbackslash ejercicio}), este entorno es para trabajo que
involucra código, datos o computación explícita.

\begin{ejerciciocomp}{Verificación de la desigualdad triangular}
\label{comp:triangular}
Implementa en Python o Julia una función que, dado un conjunto finito de
puntos en $\mathbb{R}^n$ y una función $\eqfn{d} : \eqdom{\mathbb{R}^n} \eqop{\times} \eqdom{\mathbb{R}^n}
\eqop{\to} \mathbb{R}$, verifique si $\eqfn{d}$ es una métrica en el sentido de la
\cref{def:espacio-metrico}. Prueba con:
\begin{enumerate}[label=(\alph*), leftmargin=2em]
  \item $\eqfn{d_2}(\eqk{x},\eqk{y}) = \|\eqk{x}-\eqk{y}\|_2$ (distancia euclídea, debe pasar)
  \item $\eqfn{d_\infty}(\eqk{x},\eqk{y}) = \|\eqk{x}-\eqk{y}\|_\infty$ (norma infinito, debe pasar)
  \item $\eqfn{d}(\eqk{x},\eqk{y}) = \|\eqk{x}-\eqk{y}\|_2^2$ (cuadrado de la euclídea, debe fallar)
\end{enumerate}
¿En qué caso falla el tercer candidato y con qué contraejemplo concreto?
\end{ejerciciocomp}

%───────────────────────────────────────────────────────────────────────────────
\section{El indicador de dificultad}
\label{sec:dificultad}
%───────────────────────────────────────────────────────────────────────────────

El comando \texttt{\textbackslash dificultad\{m\}\{f\}} genera un indicador
visual de dos ejes: dificultad matemática y dificultad física. Cada eje
toma un valor entre 1 y 5. Se coloca al inicio de una sección o subsección
para calibrar las expectativas del lector:

\begin{lstlisting}[style=latex, frame=none, numbers=none]
\dificultad{3}{2}   % matemática: ★★★☆☆  física: ★★☆☆☆
\end{lstlisting}

\dificultad{3}{2}

Los dos ejes son independientes. Un problema de integración numérica puede
tener baja dificultad matemática (cuadratura de Gauss es algorítmica) y
alta dificultad física (elegir el método correcto requiere intuición del
dominio). Separar los dos ejes da información más precisa que un único
número.

%───────────────────────────────────────────────────────────────────────────────
\section{Figuras y tablas dentro de entornos}
\label{sec:floats-en-entornos}
%───────────────────────────────────────────────────────────────────────────────

\begin{definicion}{Restricción de flotantes en tcolorbox}
\label{def:no-floats}
Los entornos \texttt{figure} y \texttt{table} son flotantes: invocan el
mecanismo \texttt{\textbackslash @float} de \LaTeX, que requiere estar
en \textit{outer par mode}. Dentro de una caja \texttt{tcolorbox} ese
modo no existe. El resultado es el error \texttt{Not in outer par mode}.
Los entornos afectados son: \texttt{figure}, \texttt{table},
\texttt{acadtable}, \texttt{lstlisting[float=...]}, y \texttt{longtable}.
\end{definicion}

La solución es usar los entornos no-flotantes \texttt{boxfig} y
\texttt{boxtab} definidos en \texttt{config/tcolorbox.tex}:

\begin{lstlisting}[style=latex, frame=none, numbers=none]
\begin{ejemplo}{Distribución gaussiana}
  \begin{boxfig}{fig:gauss}{Densidad $\mathcal{N}(0,1)$}
    \includegraphics[width=\linewidth]{figs/gauss.pdf}
  \end{boxfig}
  La densidad de probabilidad es $\eqfn{p}(\equ{x}) = \frac{1}{\sqrt{2\eqcon{\pi}}} e^{-\equ{x}^2/2}$.
\end{ejemplo}

\begin{definicion}{Media muestral}
  \begin{boxtab}{tab:media}{Propiedades de la media muestral}
    \begin{tabularx}{\linewidth}{l X}
      \toprule
      \headerrow
      \thead{Propiedad} & \thead{Valor} \\
      \midrule
      Sesgo   & Insesgado: $\eqop{\mathbb{E}}[\equ{\bar{X}}] = \eqk{\mu}$ \\
      Varianza & $\eqop{\mathrm{Var}}(\equ{\bar{X}}) = \eqk{\sigma}^2/\eqk{n}$ \\
      \bottomrule
    \end{tabularx}
  \end{boxtab}
\end{definicion}
\end{lstlisting}

La diferencia con los entornos flotantes es en el ancho: dentro de una
caja siempre se usa \texttt{\textbackslash linewidth} (el ancho disponible
dentro de la caja), nunca \texttt{\textbackslash textwidth}.

\begin{acadtable}[H]{tab:float-comparison}{Entornos flotantes vs.\ no-flotantes}
  \begin{tabularx}{\textwidth}{l l l X}
    \toprule
    \headerrow
    \thead{Situación} & \thead{Figuras} & \thead{Tablas} & \thead{Ancho} \\
    \midrule
    Flujo normal de texto  & \texttt{figure}    & \texttt{acadtable} & \texttt{\textbackslash textwidth} \\
    Dentro de tcolorbox    & \texttt{boxfig}    & \texttt{boxtab}    & \texttt{\textbackslash linewidth} \\
    \bottomrule
  \end{tabularx}
\end{acadtable}

%───────────────────────────────────────────────────────────────────────────────
\section{Notas bibliográficas}
\label{sec:entornos-notasbib}
%───────────────────────────────────────────────────────────────────────────────

Al final de un capítulo, el entorno \texttt{notasbib} sirve para
contextualizar históricamente los resultados presentados: atribuir
resultados a sus autores originales, indicar referencias para profundizar,
o señalar generalizaciones que no cupieron en el cuerpo principal.

\begin{lstlisting}[style=latex, frame=none, numbers=none]
\begin{notasbib}
  El teorema espectral en su forma moderna se debe a von Neumann (1929).
  Para una exposición rigurosa ver Reed y Simon [reed1980],
  Vol.~I, Cap.~VII...
\end{notasbib}
\end{lstlisting}

\begin{notasbib}
El marco general de los espacios $L^p$ fue desarrollado por Riesz
(1910) y Fischer (1907), cuya equivalencia estableció el
\textit{teorema de Riesz-Fischer}. Para una exposición moderna y rigurosa,
Reed y Simon~\cite{reed1980} es la referencia estándar en análisis
funcional orientado a la física matemática. Los espacios métricos en su
generalidad moderna se deben a Fréchet (1906); Munkres~\cite{munkres2000}
ofrece una introducción topológica accesible.
\end{notasbib}

%───────────────────────────────────────────────────────────────────────────────
\printbibliography[heading=subbibliography]
